\section{The three four-letter cores}
Take $Q=\{1,2,3,4\}$ and $I=\{0,1\}$. The three rooted cores are
\[
 A:\quad h_0=(2\ 3),\quad h_1=(1\ 2)(3\ 4),\quad d=4;
\]
\[
 B:\quad h_0=(2\ 3),\quad h_1=(1\ 2)(3\ 4),\quad d=3;
\]
\[
 C:\quad h_0=(1\ 2)(3\ 4),\quad h_1=(1\ 3)(2\ 4),\quad d=4.
\]
In each case let $\pi_S$ be the orbit partition of $H_S$ on $Q$.
For the empty set it is discrete and for $I$ it has one block.
The two singleton partitions are respectively
\[
 A,B:\quad 1\mid23\mid4,\qquad12\mid34;
 \qquad
 C:\quad12\mid34,\qquad13\mid24.
\]
Their common refinement is discrete, so the partition-meet hypothesis
of the attachment theorem in [U] holds.

In A and B the full core is the dihedral group of order eight preserving
the unordered pair $\{\{1,4\},\{2,3\}\}$.
Its intersections with the two displayed Young subgroups are
$\langle h_0\rangle$ and $\langle h_1\rangle$ respectively.
For the second assertion, among the four permutations independently
exchanging $1,2$ and $3,4$, just the identity and their simultaneous
exchange preserve that unordered pair. In C the full core is
$\{1,(12)(34),(13)(24),(14)(23)\}$, and its intersections with
the two displayed Young subgroups are again the indicated cyclic groups.
The singleton groups intersect trivially, and each nonidentity generator
fails to preserve the other singleton partition. This proves the core and mixed subgroup identities in [U], including
the empty and full sets.

For saturation, the empty-set case is immediate. A singleton generator
which is a transposition either joins $d$ to another core letter or is
disjoint from $d$; adjoining $(d\ x)$ yields exactly the appropriate
Young group. If the generator is a double transposition
$h=(u\ v)(d\ w)$, then
\[
 \bigl(h(d\ x)\bigr)^3=(u\ v),
\]
so we obtain both $(u\ v)$ and $(d\ w)$, as well as $(d\ x)$.
They generate $\Sym(\{u,v\})\times\Sym(\{d,w,x\})$.
Finally $H_I$ is transitive on $Q$. Conjugating $(d\ x)$ by its
elements gives $(q\ x)$ for every $q\in Q$, which generate
$\Sym(Q\cup\{x\})$. Thus saturation holds in every case.

\begin{corollary}
Attach any finite tree to one of these cores at its displayed root,
with that root a leaf of the attached tree. The displayed two
involutions and all its edge transpositions form a C-group marking
of $\Sym(\Omega)$ of rank $|\Omega|-2$.
\end{corollary}

\begin{proof}
The finite hypotheses were proved above, and the attached tree has
$|\Omega|-4$ edges. Apply the core-attachment theorem of [U].
\end{proof}

These are precisely the diagrams in [FL, Theorem 1.2], verified against
the accepted PDF and the authors' original TeX. Family A has the path
$1\mathbin{-}^{1}2\mathbin{-}^{0}3\mathbin{-}^{1}4$; family B
attaches at vertex $3$ instead of $4$, and C is the alternating square.
The tail has one first edge at the root, but may branch at its next
vertex. The proof gives all positive intersection assertions at once.
The following sections prove exhaustion and the complete geometric claim.

\section{Why the leaf cannot be dropped}
In core A, attach two new vertices $5$ and $6$ directly to $4$.
The marking subsets containing both core generators and just one of
$(4\ 5),(4\ 6)$ generate $\Sym(\{1,2,3,4,5\})$ and
$\Sym(\{1,2,3,4,6\})$. Their intersection is $\Sym(Q)$,
whereas the common marking subset generates only the order-eight
core. This is a literal failure of the intersection property.
The hypothesis is structural, not merely a convenience of the proof.

